% =====================================================================
%  10. Analyse comparative, limites et recommandations
% =====================================================================
\clearpage
\section{Analyse comparative et discussion}
\label{sec:discussion}

\subsection{Vue d'ensemble des approches}

Le \cref{tab:recap} rassemble, pour la limite d'élasticité, les principales
approches évaluées au cours du projet. Les autres cibles suivent le même schéma
(\cref{sec:resultats,sec:ssl}).

\begin{table}[htbp]
\centering
\small
\caption{Récapitulatif des approches pour \YS{} : $R^2$ moyen sur cinq plis externes.}
\label{tab:recap}
\begin{tabularx}{\textwidth}{@{}L c c c@{}}
\toprule
Approche & Entrées / imputation & $R^2$ aléatoire & $R^2$ groupé \\
\midrule
Référence médiane de $y$ & -- & $-0{,}031$ & $-0{,}073$ \\
Forêt aléatoire, premier benchmark & 13 / itérative & -- & 0,079 \\
Forêt aléatoire & 13 / médiane & -- & 0,102 \\
Forêt aléatoire & 16 / médiane & 0,765 & 0,390 \\
Sélection conjointe interne (schéma, imputation, famille) & 11--16 / les deux & -- & 0,337 \\
\textbf{Sélection interne de la famille (référence)} & \textbf{16 / médiane} & \textbf{0,755} & \textbf{0,337} \\
$k$NN réglé & 16 / médiane & 0,619 & $-0{,}140$ \\
COREG (semi-supervisé) & 16 / médiane & 0,629 & $-0{,}056$ \\
\bottomrule
\end{tabularx}
\end{table}

Trois enseignements en découlent :
\begin{enumerate}
  \item \textbf{Les entrées comptent davantage que l'algorithme.} Passer de 13 à 16 entrées fait passer le $R^2$ groupé de la forêt de 0,10 à 0,39, alors que changer de famille à entrées fixées modifie beaucoup moins le score.
  \item \textbf{Le protocole change la conclusion.} Toutes les familles apprises perdent entre 0,3 et 0,8 point de $R^2$ entre validation aléatoire et groupée ; les méthodes locales ($k$NN, COREG) deviennent même pires que la constante.
  \item \textbf{Le semi-supervisé n'aide pas ici.} Exploiter les entrées sans label n'apporte pas l'information qui manque pour transférer entre séries.
\end{enumerate}

\subsection{Quelle approche est la plus appropriée ?}

\begin{decision}
L'approche recommandée est la \textbf{procédure supervisée 16 entrées + procédé,
imputation médiane apprise dans le train, avec choix de la famille et des
hyperparamètres par validation interne}. En pratique, cette procédure retient
une \textbf{SVR RBF} ou une \textbf{forêt aléatoire} selon la cible
(\cref{tab:modeles_finaux}) ; la forêt est la plus robuste lorsque la série est
inconnue. COREG est conservé comme expérience semi-supervisée documentée, pas
comme modèle de référence.
\end{decision}

Le niveau de confiance à accorder dépend toutefois de l'usage :
\begin{itemize}
  \item \textbf{Soudure proche d'une série déjà documentée} (même procédé, mêmes familles de composition) : l'erreur typique attendue est d'environ 30\,MPa sur \YS, 25\,MPa sur \UTS{} et 2 points d'allongement --- un ordre de grandeur utile pour pré-sélectionner des recettes avant essais.
  \item \textbf{Soudure d'une nouvelle étude, procédé ou famille d'alliage} : l'erreur double (environ 55\,MPa sur \YS{} et \UTS) et le $R^2$ chute sous 0,35. Le modèle ne remplace pas l'essai mécanique ; il peut au mieux orienter.
  \item \textbf{Énergie Charpy} : les prédictions restent exploratoires tant que les protocoles ne sont pas harmonisés ; elles ne doivent pas servir à une décision de ténacité.
\end{itemize}

\subsection{Limites de l'étude}

\begin{itemize}
  \item \textbf{Sélection rétrospective.} Le passage à 16/médiane a été décidé après observation des scores groupés sur ces mêmes données. Même avec une validation imbriquée ensuite, l'ensemble de la démarche n'a pas été validé sur un jeu indépendant.
  \item \textbf{Groupes imparfaits.} Les séries sont inférées des identifiants : elles peuvent fusionner ou séparer des publications. Elles restent des approximations des vraies unités d'indépendance.
  \item \textbf{Données hétérogènes.} Compositions incomplètes, unités douteuses en quarantaine, protocoles Charpy non harmonisés, précision de transcription inconnue.
  \item \textbf{Entrées omises.} Des éléments exclus pour leur faible couverture (Nb, V, Ti, B, Cu\ldots) peuvent expliquer une partie des décalages entre séries.
  \item \textbf{Incertitude.} Les intervalles rapportés sont descriptifs ; aucune prédiction n'est accompagnée d'un intervalle de prédiction calibré.
  \item \textbf{Semi-supervisé.} Une seule méthode, un seul budget de labels (complet) ; d'autres variantes (régresseurs différents, labels rares) restent à explorer.
\end{itemize}

\subsection{Recommandations}
\label{subsec:recommandations}

L'énoncé demande nos recommandations pour obtenir une bonne qualité de soudure.
Les modèles mettent en évidence des \emph{associations}, pas des causes : nous
séparons donc ce que les données permettent d'affirmer de ce qui demanderait des
essais dédiés.

\paragraph{Ce que l'étude permet de recommander.}
\begin{enumerate}
  \item \textbf{Mesurer et enregistrer la composition complète}, y compris Ni, Cr et Mo : leur ajout réduit l'erreur de 9 à 11\,MPa sur la résistance, alors qu'ils manquent dans plus de la moitié des lignes. Une valeur mesurée vaut mieux qu'une valeur imputée.
  \item \textbf{Documenter le traitement thermique après soudage} (température, durée) : avec la composition, c'est le bloc d'information le plus utilisé par les modèles pour prédire la résistance.
  \item \textbf{Enregistrer les protocoles d'essai}, en particulier pour Charpy : distinguer une énergie mesurée à température imposée d'une température de transition, et conserver des courbes $E(T)$ complètes.
  \item \textbf{Identifier les séries, éprouvettes et publications} dans les bases, afin de pouvoir valider honnêtement le transfert à de nouvelles études.
  \item \textbf{Valider tout modèle sur de nouvelles séries} avec de vraies mesures avant un usage opérationnel, et signaler les soudures hors du domaine du train (procédé inconnu, compositions extrêmes), où l'erreur augmente nettement.
  \item \textbf{Privilégier la diversité des études} plutôt que le nombre de lignes : les courbes d'apprentissage suggèrent que de nouvelles séries sont plus informatives que de nouvelles répétitions.
\end{enumerate}

\paragraph{Ce que l'étude ne permet pas d'affirmer.} Aucun réglage de soudage
\og optimal \fg{}, aucun seuil universel de bonne soudure et aucun effet causal
d'un élément ou d'un paramètre n'est démontré. Les importances indiquent
\emph{quelles} informations le modèle utilise, pas \emph{dans quel sens} modifier
une recette. Pour formuler des recommandations de procédé (composition cible,
apport thermique, PWHT), il faudrait des plans d'expériences contrôlés et
l'expertise de métallurgistes du soudage.

\begin{aretenir}
La procédure supervisée 16/médiane (SVR ou forêt aléatoire selon la cible) est
l'approche la plus appropriée. Elle est fiable pour interpoler dans des séries
connues, nettement moins pour de nouvelles études. Les leviers d'amélioration
sont avant tout la \textbf{qualité et la complétude des données} : composition
complète, PWHT, protocoles d'essai et identification des séries.
\end{aretenir}
